\newpage
\section{REQUISITOS}
\label{sec:requisitos}
Os requisitos apresentados buscam apresentar as necessidades exergadas pelo grupo em relação ao uso da plataforma por pessoas idosas e outros agentes relacionados. Os requisitos foram dividos em requisitos funcionais e não funcionais.

\subsection{Requisitos Funcionais}
Os requisitos funcionais definem os recuros que o sistema deve fornecer. Sendo eles:

\begin{itemize}
    \item \textbf{Acesso e Autenticação}
        \begin{enumerate}
            \item [US1] Como pessoa idosa com pouca experiência digital, quero acessar a plataforma em até 2 passos após abrir o aplicativo, para não me sentir confuso ou desistir do uso.
            \item [US2] Como usuário, quero recuperar minha senha de forma simples, para não ficar bloqueado fora da plataforma.
            \item [US3] Como idoso, quero acessar a plataforma com login simples (ex: botão único ou biometria), para evitar dificuldades com senhas complexas.
        \end{enumerate}
 
    \item \textbf{Visualização e Inscrição em Atividades}
        \begin{enumerate}
            \item [RF1] O sistema deve permitir que o participante visualize atividades disponíveis.
            \item [RF2] O sistema deve permitir inscrição em atividades.
        \end{enumerate}
 
    \item \textbf{Participação em Atividades}
        \begin{enumerate}
            \item [US4] Como pessoa idosa, quero participar de atividades online com apenas um clique a partir da tela inicial, para facilitar meu engajamento e evitar dificuldades técnicas.
            \item [US5] Como pessoa idosa, quero participar de atividades online (exercícios, jogos, encontros), para manter minha saúde e interação social.
            \item [RF3] O sistema deve permitir acesso a atividades online (ex: videochamada).
        \end{enumerate}
 
    \item \textbf{Videochamadas e Comunicação}
        \begin{enumerate}
            \item [US6] Como pessoa idosa, quero entrar em chamadas de vídeo automaticamente ao clicar em um botão, sem necessidade de configurações adicionais, para interagir com outras pessoas sem dificuldades técnicas.
            \item [US7] Como pessoa idosa, quero enviar mensagens de voz ou texto para outros participantes, para me comunicar de forma simples e manter vínculos sociais.
            \item [RF4] O sistema deve permitir troca de mensagens entre participantes.
            \item [RF5] O sistema deve permitir envio de mensagens em grupo.
        \end{enumerate}
 
    \item \textbf{Lembretes e Notificações}
        \begin{enumerate}
            \item [RF6] O sistema deve enviar lembretes automáticos de atividades (notificações ou áudio), para que os usuários não esqueçam dos eventos programados.
        \end{enumerate}
 
    \item \textbf{Conteúdos e Materiais}
        \begin{enumerate}
            \item [US8] Como idoso, quero acessar gravações das atividades passadas, para participar mesmo quando não puder estar presente ao vivo.
            \item [RF7] O sistema deve permitir acesso a conteúdos compartilhados.
            \item [RF8] O sistema deve permitir upload de conteúdos.
            \item [US9] Como pessoa idosa, quero acessar tutoriais simples e guiados, para aprender a usar a plataforma sem dificuldades.
        \end{enumerate}
 
    \item \textbf{Criação e Gestão de Atividades (Mediador/Tutor)}
        \begin{enumerate}
            \item [US10] Como mediador, quero criar e agendar atividades com data, horário e descrição, para organizar a participação dos membros de forma clara.
            \item [RF9] O sistema deve permitir criação e edição de atividades.
            \item [US11] Como tutor, quero poder marcar em calendário eventos comunitários com avisos, para que os membros estejam cientes do que ocorrerá.
        \end{enumerate}
 
    \item \textbf{Monitoramento e Engajamento (Mediador/Tutor)}
        \begin{enumerate}
            \item [US12] Como mediador, quero visualizar quem participou das atividades, para acompanhar o engajamento dos membros.
            \item [RF10] O sistema deve fornecer relatórios de participação.
            \item [US13] Como mediador, quero enviar mensagens motivacionais e avisos para o grupo, para estimular a participação contínua.
            \item [RF11] O sistema deve permitir gerenciar lista de participantes.
        \end{enumerate}
 
    \item \textbf{Moderação}
        \begin{enumerate}
            \item [US14] Como mediador, quero moderar interações no chat ou nas atividades, para garantir um ambiente seguro e respeitoso.
        \end{enumerate}
 
    \item \textbf{Social e Comunidade}
        \begin{enumerate}
            \item [US15] Como sênior, quero conhecer membros com gostos parecidos, para criar novas amizades.
            \item [US16] Como sênior, quero compartilhar acontecimentos do meu dia para outros membros, para que outros vejam o que estou fazendo.
            \item [US17] Como sênior, quero poder organizar eventos presenciais pela plataforma, para poder reunir membros da comunidade.
            \item [US18] Como sênior, quero poder conhecer membros que estejam localmente próximos de mim, para poder facilitar a criação de laços comunitários.
            \item [US19] Como tutor, quero poder selecionar para que membros minhas atividades são compartilhadas, para poder reunir membros com gostos semelhantes.
        \end{enumerate}
 
    \item \textbf{Gestão de Usuários (Administração)}
        \begin{enumerate}
            \item [US20] Como tutor, quero poder alterar os dados dos membros para atualizar informações dos usuários.
            \item [RF12] O sistema deve permitir cadastro e remoção de usuários.
            \item [RF13] O sistema deve controlar permissões de acesso.
            \item [RF14] O sistema deve registrar logs de atividades.
        \end{enumerate}
 
    \item \textbf{Suporte ao Usuário}
        \begin{enumerate}
            \item [US21] Como idoso, quero ter acesso a um botão de ajuda ou suporte imediato, para resolver dúvidas ou dificuldades rapidamente.
        \end{enumerate}
 
    \item \textbf{Preferências e Personalização}
        \begin{enumerate}
            \item [US22] Como idoso, quero configurar preferências (volume, tamanho da fonte, notificações), para adaptar a plataforma às minhas necessidades.
        \end{enumerate}
\end{itemize}

\subsection{Requisitos Não Funcionais}
Os requisitos não funcionais focam em atributos que o sistema deve aparesentar. Sendo eles:

\begin{itemize}
    \item \textbf{Usabilidade}
        \begin{enumerate}
            \item [RNF1] A interface deve ser simples e intuitiva, com navegação clara e direta.
            \item [RNF2] Deve-se minimizar o número de passos para ações principais.
        \end{enumerate}
 
    \item \textbf{Acessibilidade}
        \begin{enumerate}
            \item [RNF3] Deve usar botões grandes e textos legíveis (mínimo de 16px).
            \item [RNF4] Suporte a alto contraste.
            \item [RNF5] Possibilidade de aumentar fonte.
            \item [RNF6] Compatível com leitores de tela.
        \end{enumerate}
 
    \item \textbf{Desempenho}
        \begin{enumerate}
            \item [RNF7] O sistema deve carregar rapidamente mesmo em conexões lentas.
        \end{enumerate}
 
    \item \textbf{Segurança e Privacidade}
        \begin{enumerate}
            \item [RNF8] Dados dos usuários devem ser protegidos e sua privacidade garantida.
            \item [RNF9] Autenticação segura.
            \item [US23] Como idoso, quero interagir em um ambiente moderado e seguro, para evitar golpes, spam ou conteúdos inadequados.
        \end{enumerate}
 
    \item \textbf{Compatibilidade}
        \begin{enumerate}
            \item [RNF10] Deve funcionar em celulares, tablets e computadores.
        \end{enumerate}
\end{itemize}